\subsection{A finite boundary kernel for simultaneous whole-cut cancellation}
\label{marked-cut:boundary-kernel}

The separate primitive-family tests can miss a sequence in which cancelling
one annular gluing exposes a useful disk for another. We now put all original
core two-handles and all annular gluing two-handles into one geometric
cancellation problem. The bundle genus can be exponentially large, but only
one column per named boundary circle needs to be retained. The construction
below supplies the geometric justification missing from the earlier
boundary-template diagnostic.

\begin{proposition}[Bordered-bundle boundary disk system]
\label{marked-cut:bundle-disk-system}
Let $B$ be an orientable interval bundle over a connected compact surface
$F$ with nonempty boundary, $b=|\partial F|$, and $G=1-\chi(F)$.
Let $a\le b$ distinct vertical boundary annuli be named. There is a complete
system of $G$ meridian disks with the following unsigned intersections with
the annular midpoint curves.
\begin{enumerate}
\item If $a<b$, the first $a$ named curves each meet their own disk once;
all other intersections of named curves with the system vanish.
\item If $a=b$, choose the last named boundary as distinguished. Each of
the first $a-1$ curves meets its own disk once. The distinguished curve
meets each of those disks once and each of the other $G-(a-1)$ disks twice.
All remaining intersections vanish.
\end{enumerate}
The disks can avoid every specified proper interval-attachment region.
This includes nonorientable bases whose interval bundle has orientable total
space; no product-bundle assumption is required.
\end{proposition}
\begin{proof}
Choose a distinguished boundary $C_0$, unnamed when $a<b$. Cut $F$ along
disjoint arcs from each of its other boundary components to $C_0$. The
result is a one-boundary surface. A standard system of $2g$ further arcs
for an orientable genus-$g$ surface, or $q$ further arcs for a nonorientable
surface of genus $q$, cuts it to a disk. All endpoints of these latter arcs
lie on $C_0$. Equivalently, take a disk with bands and $b-1$ boundary holes;
the band cocores and arcs joining the holes to its outer boundary are such
a system. This gives exactly $G$ proper arcs cutting $F$ to a disk.

The interval preimage of each arc is a disk, even in the twisted bundle.
Their complement is the interval bundle over a disk and hence a ball.
An endpoint on a boundary circle contributes exactly one intersection with
its vertical midpoint curve. When all boundaries are named, the $b-1$
joining arcs have one endpoint on $C_0$, and each remaining arc has two.
When $a<b$, choose $C_0$ unnamed and order the joining arcs with named
boundaries first; unnamed-boundary and genus arcs have no endpoints on
named circles. This gives the stated matrices geometrically, without using
free reductions of surface-group words.

Each boundary circle that is not fully attached has a genuine open gap.
Move the finitely many required endpoints into such gaps. Their order can
be chosen inside a gap, so all arcs stay disjoint and their vertical disks
avoid proper attachment intervals. If every boundary is fully attached,
there are no proper interval attachments on that base.
\end{proof}

\begin{proposition}[Source annulus gluing has a finite projected handle model]
\label{marked-cut:projected-gluing}
Assume the authenticated source core has no three-handles and all bundle
bases have nonempty boundary. Contract a spanning forest of its zero/one
handles. For every ported bundle add one zero-handle and the disk system
of Proposition~\ref{marked-cut:bundle-disk-system}. For each proper disk or
full annulus interface add a connecting one-handle; for each full annulus
also add one two-handle. This is a handle decomposition of the actual cut.
After a forest of connecting handles is contracted, the retained columns
are the original surviving core one-handles and at most $P$ named bundle
boundary columns. Every omitted column has initial entries only zero or two.
\end{proposition}
\begin{proof}
The source theorem provides embedded, mutually disjoint interface disks and
annuli, disjoint from the original two-handle attaching regions. A disk
identification is the usual boundary connecting one-handle description.
For an annulus, take a narrow spanning rectangle in each copy, attach the
connecting handle there, and identify the remaining rectangles by one
two-handle. Its attaching curve follows the two annular midpoint curves,
with the small spanning rectangles removed, and crosses the connecting
handle twice. The rectangles may be chosen between consecutive disk-system
crossings. Thus its intersections with the old core disks are exactly the
source midpoint crossings, and its intersections with bundle disks are
exactly the preceding proposition's counts. The real core two-handles have
zero intersections with bundle and connecting cocores. All these operations
are simultaneous on disjoint source interfaces.

A forest of connecting handles merges balls. Removing its cocores simply
deletes those columns; no other geometric intersection count changes.
Each remaining connecting column has a two in its own annular row, or is
zero throughout for a proper disk interface. The omitted bundle columns
are zero on all rows if $a<b$, and have a two only in the distinguished row
if $a=b$. A bundle with no annuli contributes only omitted zero columns.
Consequently, regardless of the binary size of $G$, the retained matrix
has $O(T+P)$ rows and columns and is exactly the projection of a genuine
geometric handle matrix for the supplied cut.
\end{proof}

\begin{proposition}[Safe omission and a handlebody census]
\label{marked-cut:kernel-census}
Apply only retained singleton pivots, with saturated update
$n'_{ih}=\min(2,n_{ih}+n_{ig}n_{rh})$ after pivot $(r,g)$.
Every such pivot lifts to a geometric cancellation in the complete handle
model. If all two-handle rows disappear, every component of the cut is a
handlebody. Its genus is the number of surviving retained columns plus its
omitted bundle genus and its connecting-graph cycle rank.
\end{proposition}
\begin{proof}
The geometric disk-slide lemma applies to the complete matrix. Projection
onto retained columns commutes with its positive update: a retained entry
uses only other retained entries and the retained pivot column. Thus a
singleton seen by the finite replay is an actual unique belt intersection.
Omitted entries remain in $\{0,2\}$ after saturation, since an update only
adds nonnegative multiples of entries in that set. They never become
singleton candidates. They may grow, and their disks may be slid, but they
are never deleted. No exponentially long word or disk-slide list is needed
to verify the retained incidence predicates.

When all rows have cancelled there are no two- or three-handles. Forest
contraction left one zero-handle per connected attachment-graph component.
The remaining orientable zero/one-handle manifold is a handlebody. Counting
its surviving handles gives the formula. Equivalently, for a connected
graph piece with vertices $V$ and interfaces $E$,
\[
 g=\sum_{c}g_c^{(01)}+\sum_F(1-\chi(F))
       +|E|-|V|+1-h_2-a,
\]
where $h_2$ counts the original core two-handles and $a$ its full annular
interfaces. The two counts agree because every verified pivot removes one
retained one-handle and one actual two-handle. Isolated bundles contribute
binary multiplicities separately. The original source Euler identity gives
an additional consistency check, not the proof of handlebody topology.
\end{proof}

\paragraph{Independent replay.}
The whole-cut schema \path{normal-cut-handlebodies-v5} contains the two
source certificates, a \path{bundle-boundary-kernel-v1} record, and the
census. The kernel records both forests, every annular source crossing,
the named boundary columns and omitted genus counts, and the pivot schedule.
The initial v5 producer used sparse rows and tried projected row-leaf
peeling before a positive-fill schedule. Section~\ref{marked-cut:column-peeling}
replaces that search with complete column peeling. A projected leaf can still meet omitted
columns twice; the checker therefore always uses the general positive
update, without assuming that the full row has only one crossing.

The independent module \path{bundle_boundary_kernel_verify} reconstructs
source midpoint paths, derives the canonical surface templates, builds its
own forests, and replays a dense matrix. It imports no kernel producer,
cancellation planner, or census helper. The producer counts surviving
columns; the checker also counts initial handles minus all two-handles,
component by component, and checks agreement. Existing successful v1--v4
certificates remain preferred and unchanged; v5 is attempted when their
core cancellation or complete private-family test fails.

There are $O(T+P)$ retained handles, so at most $O(T+P)$ pivots need
$O((T+P)^3)$ saturated-entry operations. This is in addition to reading the source-query output and the binary
arithmetic for genus and multiplicity. The retained matrix is quadratic
in size; source crossings and template records have polynomial size.
Omitted genus and isolated multiplicity use exact binary arithmetic, never
iteration over their values. These bounds are for a supplied normal cut.
The raw full-row decision is now complete for the supplied matrix, but
three-handles are still unsupported,
and no general bound on useful-surface discovery or hierarchy length follows.

\paragraph{Source audit and focused validation.}
The two cases identified by the earlier diagnostic now have source-certified
v5 schedules and independently replayed censuses. The nine-tetrahedron mixed
Fibonacci source yields one ball and two genus-three handlebodies; the
ten-tetrahedron mixed source yields one ball and two genus-two handlebodies.
The expanded-column regression restores the omitted genus and connecting
columns on these finite examples and checks every pivot with unsaturated
nonnegative counts. All 37 focused tests pass in 73.438 seconds, including hostile forest,
curve, template, schedule and census records; hexadecimal transport;
disabled producers; source relabellings; and binary multiplicity $2^{1024}$.

The 106-case source audit now accepts 85 cuts, with independent replay for
all 85. Literal Regina cuts agree on all 78 accepted cases below the
4,096-disk expansion cap; seven large accepted inputs are intentionally not
expanded. All 99 explicit oracle runs complete without timeout. The 83
previously accepted certificates have identical canonical hashes. The 21
remaining refusals comprise eight stalled raw schedules, three square raw
obstructions, and ten cores with three-handles. No previous success is lost.
The source-pinned records are retained in \path{data/boundary-kernel-audit.json}
and the focused log in \path{data/boundary-kernel-focused.txt}.
The full Regina-enabled regression passes all 1,678 tests, with no skips,
in 570.513 seconds. Its 13,846 expanded-sheet comparisons and 196,712
marked-cover comparisons also pass. The full log is retained in
\path{data/boundary-kernel-tests.txt}. No performance comparison is inferred
from this concurrently executed validation audit.

A separate direct-kernel audit exercises v5 on all 96 sources without
three-handles, including those for which production prefers v1--v4. It
constructs 85 source certificates, independently replays all of them and
matches the 78 available literal oracle histograms. The other eleven
schedules conservatively fail. The direct certificates need not equal the
preferred older certificates; the ordinary producer still preserves those
older records. The driver and pinned results are retained in
\path{data/boundary-kernel-direct-audit.py} and its companion JSON.

\paragraph{Paired timings and binary capacity.}
After regression, direct audits and the implementation article build
finished, six finite cuts were measured in five shuffled paired rounds after
warmup. Native time includes whole-cut certificate production and independent
replay; the explicit path constructs Regina's actual cut and recognizes each
component. Parsing and serialization are excluded. Every histogram agrees.
Three identical-code explicit A/A pairs accompany every case. The driver and
all samples are retained in \path{data/boundary-kernel-benchmark.py} and its
companion JSON, including the corpus, driver and runtime source hashes.

\begin{table}[htbp]
\centering
\begin{tabular}{rrrrrr}
\toprule
$T$ & Copies & Native (s) & Explicit (s) & Ratio & A/A \\
\midrule
9 & 1 & 2.3270 & 0.0416 & 0.018 & 0.968 \\
9 & 16 & 10.1010 & 1.3882 & 0.136 & 1.003 \\
9 & 256 & 10.3384 & 23.5798 & 2.281 & 1.003 \\
10 & 1 & 7.5968 & 0.0407 & 0.007 & 1.026 \\
10 & 16 & 10.1140 & 0.4376 & 0.028 & 0.973 \\
10 & 256 & 12.4218 & 5.3981 & 0.448 & 0.975 \\
\bottomrule
\end{tabular}
\caption{Whole-cut v5 queries for the two newly certified mixed families. Native includes production and independent replay. Ratio is the median of paired explicit/native times; A/A compares identical explicit cuts.}
\label{marked-cut:boundary-kernel-timings}
\end{table}


The nine-tetrahedron source at 256 copies has median native time 10.338
seconds and explicit time 23.580 seconds, with paired ratio 2.281. The other
five cases favor explicit construction, with paired ratios from 0.007 to
0.448. A/A medians range from 0.968 to 1.026, but native timing also varies:
for example, the ten-tetrahedron one-copy samples span 4.568--8.329 seconds,
and its 256-copy samples span 10.200--17.691 seconds. The table reports medians
of paired ratios, which need not equal the ratios of the two marginal medians.
These measurements support a gain for one supplied family and regressions
for the others; they are not a universal speedup or full unknot-recognition
benchmark.

For each family a separate single-run capacity check uses multiplier
$N=2^{16384}$. The nine-tetrahedron query takes 46.349 seconds for production
plus replay and its serialized certificate contains 138,607,389 bytes.
Its census is $N$ balls, two genus-three handlebodies, and $N-1$ genus-five
handlebodies. The ten-tetrahedron query takes 43.830 seconds and serializes
to 137,651,867 bytes; it certifies $N$ balls, two genus-two handlebodies, and
$N-1$ genus-three handlebodies. No literal expansion is attempted.

The two added kernels have respectively 24 and 27 two-handle pivots,
eight and nine annular rows, five and three retained boundary columns,
and one and zero omitted genus columns. Thus the binary multiplicity does
not require an expanded handle matrix. The large complete certificates also
contain all source-orbit proofs; their roughly 138 MB serialized sizes are
a practical limitation that the small final cancellation record does not
remove. The publication consistency script checks the exact binary censuses,
all runtime pins, previous certificate hashes, and the retained validation
artifacts in \path{data/boundary-kernel-publication-check.json}.

\paragraph{Next obstruction: a source two/three-handle diagnostic.}
For the ten remaining sources with three-handles, a separate probe assigns
ends to each extreme source edge gap. With zero normal weight, its two ends
are the original edge endpoints; with positive weight, one end is the
corresponding original vertex and the other lies on the cut boundary.
The probe repeatedly deletes a candidate two/three pair when exactly one
end belongs to a remaining interior vertex handle, counting repeated ends
with multiplicity. It then tests the remaining rows with the boundary kernel.

All ten diagnostic reductions remove every three-handle. Seven resulting
kernels eliminate all two-handles and match the literal handlebody genera.
The other three kernels fail, and each corresponding literal cut contains a
non-handlebody component. These are promising candidate reductions, not seven
additional accepted cuts: independent authentication of the dual endpoints
and replay of the geometric two/three moves remain to be implemented. The
production count remains 85. The diagnostic driver and records are retained
in \path{data/core-three-handle-probe.py} and its companion JSON.

